\subsection{Definition of the coequalizer}

Let H be a quiver, let G be a groupoid, and let \(\varphi\) , \(\psi\) be morphisms of quivers from H to G. Denote by \(\mathrm{Coh}(\varphi,\psi)\) the cohomotoper of the pair \((\varphi,\psi)\) (II, p. 185, def. 3), \(\alpha\colon\mathrm{G}\to\mathrm{Coh}(\varphi,\psi)\) the canonical groupoid morphism and \(h\) the canonical homotopy relating \(\alpha\circ\varphi\) to \(\alpha\circ\psi\) .

Let \(\operatorname{Coeg}(\varphi,\psi)\) be the groupoid obtained from \(\operatorname{Coh}(\varphi,\psi)\) by contracting the arrows belonging to the image of h (II, p. 196, no. 1); let \(\beta\colon\operatorname{Coh}(\varphi,\psi)\to\operatorname{Coeg}(\varphi,\psi)\) denote the canonical morphism and set \(\gamma=\beta\circ\alpha\) .

DEFINITION 2. --- We say that the groupoid \(\operatorname{Coeg}(\varphi,\psi)\) is the coequalizer of the pair \((\varphi,\psi)\); the groupoid morphism \(\gamma\) is called the canonical morphism from G to \(\operatorname{Coeg}(\varphi,\psi)\).

PROPOSITION 3. --- The pair (Coeg(φ,ψ),γ) possesses the following universal property:

\begin{enumerate}
\def\labelenumi{\alph{enumi})}
\item
  We have \(\gamma \circ \varphi = \gamma \circ \psi\) .
\item
  Let G' be a groupoid and \(\theta\colon G\to G'\) a morphism of groupoids such that \(\theta\circ\varphi=\theta\circ\psi\) . There exists a unique groupoid morphism \(\overline{\theta}\colon\operatorname{Coeg}(\varphi,\psi)\to G'\) such that \(\overline{\theta}\circ\gamma=\theta\) .
\end{enumerate}

Let \(a\) be a vertex of H. The arrow \(h(a)\) of \(\mathrm{Coh}(\varphi, \psi)\) connects \(\alpha(\varphi(a))\) to \(\alpha(\psi(a))\) . By definition of the groupoid \(\mathrm{Coeg}(\varphi, \psi)\) , the origin and the target of the arrow \(\beta(h(a))\) are equal; we therefore have \(\gamma(\varphi(a)) = \gamma(\psi(a))\) .

Let f be an arrow of H; if a denotes its origin and b its target, we thus have

\[
\alpha (\varphi (f)) \cdot h (b) = h (a) \cdot \alpha (\psi (b)).
\]

Taking the image under \(\beta\) of both sides of this equality, we obtain the relation \(\gamma(\varphi(f)) = \gamma(\psi(f))\) . This proves the assertion \(a\) ).

Let us prove \(b\) ). The map \(\eta\colon\mathrm{Som}(\mathrm{H})\to\mathrm{Fl}(\mathrm{G}')\) which associates to every vertex \(a\) of H the arrow \(e_{\theta(\varphi(a))}\) of G' is a homotopy connecting \(\theta\circ\varphi\) to \(\theta\circ\psi\). By the universal property of cohomotopers (II, p. 185, prop. 3), there exists a unique morphism of groupoids \(\theta_1\colon\mathrm{Coh}(\varphi,\psi)\to\mathrm{G}'\) such that \(\theta_1\circ\alpha=\theta\) and \(\theta_1(h(a))=e_{\theta(\varphi(a))}\) for every vertex \(a\) of H. By prop. 7 of II, p. 176, this last property implies the existence of a unique morphism of groupoids \(\overline{\theta}\colon\mathrm{Coeg}(\varphi,\psi)\to\mathrm{G}'\) such that \(\overline{\theta}\circ\beta=\theta_1\). One then has \(\overline{\theta}\circ\gamma=\theta_1\circ\alpha=\theta\).

Conversely, if \(\overline{\theta}'\colon \mathrm{Coeg}(\varphi,\psi)\to \mathrm{G}'\) is a morphism of groupoids such that \(\overline{\theta}'\circ\gamma=\theta\) , we have \((\overline{\theta}'\circ\beta)\circ\alpha=(\overline{\theta}\circ\beta)\circ\alpha\) , whence \(\overline{\theta}'\circ\beta=\overline{\theta}\circ\beta\) by II, p. 185, prop. 3, whence \(\overline{\theta}'=\overline{\theta}\) by prop. 7 of II, p. 176. This proves the uniqueness of \(\overline{\theta}\) , whence the assertion \(b\) ).

COROLLARY. --- The groupoid Coeg(φ,ψ) is generated by the image of the morphism γ.

Let C be the subgroupoid of Coeg( \(\varphi,\psi\) ) generated by the image of G; denote by i the canonical morphism from C to Coeg( \(\varphi,\psi\) ) and \(\theta\colon G\to C\) the morphism such that \(i\circ\theta=\gamma\) . By Proposition 2, there exists a unique groupoid morphism \(\overline{\theta}\colon\) Coeg( \(\varphi,\psi\) ) \(\to\) C such that \(\overline{\theta}\circ\gamma=\theta\) . Then, \(\gamma=i\circ\overline{\theta}\circ\gamma\) , therefore \(i\circ\overline{\theta}\) is the identity morphism of Coeg( \(\varphi,\psi\) ) (loc. cit.). In particular, the maps Som(i) and Fl(i) are surjective, hence C = Coeg( \(\varphi,\psi\) ).

Remarks. --- 1) The set of vertices of Coeg( \(\varphi,\psi\) ) is the set of connected components of the quiver

\[
\Gamma = (\operatorname{Som} (\mathrm{G}), \operatorname{Som} (\mathrm{H}), \operatorname{Som} (\varphi), \operatorname{Som} (\psi))
\]

(II, p. 197). In other words, it is the quotient set of \(\operatorname{Som}(\mathbf{G})\) by the finest equivalence relation such that \(\varphi(a)\) is equivalent to \(\psi(a)\) for all \(a \in \operatorname{Som}(\mathbf{G})\) .

\begin{enumerate}
\def\labelenumi{\arabic{enumi})}
\setcounter{enumi}{1}
\tightlist
\item
  The map \(\mathrm{Orb}(\gamma)\) , deduced from \(\gamma\) by passing to orbits, defines by passage to the quotient a bijection from the set of connected components of the frame of the pair \((\varphi,\psi)\) to the set of orbits of the coequalizer. This follows from II, p. 185, prop. 4 and from the fact that the map deduced from \(\beta\) by passage to orbits is bijective (II, p. 170, remark 1).
\end{enumerate}

Consequently, the map from Som(G) to \(\mathrm{Orb}(\mathrm{Coeg}(\varphi,\psi))\) deduced from \(\gamma\) identifies the set of orbits of \(\mathrm{Coeg}(\varphi,\psi)\) with the quotient set of Som(G) by the equivalence relation generated by the pairs \((\varphi(x),\psi(x))\) for \(x\in\mathrm{Som(H)}\) and the pairs \((o(f),t(f))\) for \(f\in\mathrm{Fl(G)}\) .

\begin{enumerate}
\def\labelenumi{\arabic{enumi})}
\setcounter{enumi}{2}
\tightlist
\item
  The map Fl(γ) from Fl(G) to Fl(Coeg(φ,ψ)) is not in general surjective.
\end{enumerate}

